\part{Kvantová mechanika (nerelativistická teorie)}
\localtableofcontents

\chapter{Základní koncepty kvantové mechaniky}
\section{Princip neurčitosti}
Pokusíme-li se vysvětlit atomové jevy prostřednictvím klasické mechaniky a elektrodynamiky, získáváme výsledky v očividném nesouladu s experimentem. Jeden z nich obdržíme aplikací obyčejné elektrodynamiky na model atomu, v němž se elektrony pohybují okolo jádra na klasických obězných drahách. Při takovémto pohybu by měly elektrony spojitě vyzařovat elektromagnetické vlny, jakožto při jakémkoliv zrychleném pohybu. Tímto by však elektrony postupně ztratily svou energii a spadly do jádra. Dle klasické elektrodynamiky je tedy atom nestabilní, což je spor s realitou.\par
Tento rozpor mezi teorií a experimentem ukazuje, že konstrukce teorie aplikovatelné na atomové jevy--to jest jevy částic velmi malých hmotností na velmi malých vzdálenostech--nutně vyžaduje fundamentální modifikace základních fyzikálních konceptů a zákonů.\par
Jakožto výchozí bod pro zkoumání těchto modifikací je vhodné zvolit experimentálně pozorovaný fenomén zvaný \emph{elektronová difrakce}.\footnote{Fenomén elektronové difrakce byl ve skutečnosti objeven až po vynálezu kvantové mechaniky. Pro naši diskusi však není klíčová historická posloupnost vývoje teorie, nýbrž snaha ji konstruovat takovým způsobem, aby byly patrné souvislosti mezi základními principy kvantové mechaniky a experimentálními pozorováními.} Je pozorováno, že při průchodu homogenního svazku elektronů krystalem dochází ve výstupním svazku ke vzniku střídajících se maxim a minim intenzity, velmi podobně jako při difrakci elektromagnetických vln. Za určitých podmínek tedy chování hmotných částic--zde elektronů--vykazuje vlastnosti typické pro vlnové procesy.\par
Absurdita tohoto fenoménu v kontextu tradičních představ o pohybu je nejpatrnější v následujícím myšlenkovém experimentu, jež je idealizací elektronové difrakce v krystalu. Představme si stínítko, které je až na dvě vyříznuté štěrbiny pro elektrony nepropustné. Při průchodu svazku elektronů\footnote{Předpokládáme tak řídký svazek, že interakce částic v rámci něj nehraje žádnou roli.} jednou ze štěrbin, kdy druhá štěrbina je zakrytá, pozorujeme na stínítku umístěném za štěrbinou určité rozložení intenzity; stejně tak pozorujeme obecně jiné rozložení intenzite, zakryjeme-li první ze štěrbin a odkryjeme-li druhou. Při průchodu svazku oběma štěrbinami bychom na základě klasických představ předpokládali rozložení intenzity ve formě součtu předchozích dvou: každý elektron projde jednou ze štěrbin a nemá žádný vliv na elektrony procházející štěrbinou druhou. Fenomén difrakce však ukazuje, že ve skutečnosti obdržíme difrakční obrazec, jež vlivem interference naprosto neodpovídá součtu obrazců obdržených při průchodu svazků pouze jednou ze štěrbin. Tento výsledek je patrně neslučitelný s představou, že se elektrony pohybují po drahách.\par
Mechanika popisující atomové jevy--\emph{kvantová mechanika} nebo \emph{vlnová mechanika}--tudíž musí být založena na představách o pohybu, které jsou fundamentálně odlišné od představ klasických. V kvantové mechanice neexistuje koncept trajektorie částice. Tento výrok tvoří obsah tzv. \emph{principu neurčitosti}, jednoho ze základních principů kvantové mechaniky, který byl objeven W. Heisenbergem roku 1927.\footnote{Je zajímavé zmínit, že úplný matematický formalismus kvantové mechaniky byl zkonstruován W. Heisenbergem a E. Schrödingerem v letech 1925--6, tedy před objevem principu neurčitosti, který pouze objasnil fyzikální obsah tohoto formalismu.}\par
Zamítnutím obvyklých představ klasické mechaniky princip neurčitosti nabývá záporného obsahu. Samotný tedy nepostačuje jako základ pro konstrukci nové mechaniky částic. Taková teorie musí být přirozeně založena na nějakých positivních předpokladech, jež budeme diskutovat níže (§2). Pro formulaci těchto předpokladů musíme však nejdříve upřesnit znění úloh, se kterými se kvantová mechanika potýká. Proto se první věnujme speciálnímu vzájemnému vztahu kvantové a klasické mechaniky. Obecnější teorii lze obvykle formulovat logicky kompletním způsobem nezávisle na méně obecné teorii, která tvoří její limitní případ. Například lze zkonstruovat relativistickou mechaniku na základě vlastních fundamentálních předpokladů, a to bez jakéhokoliv odkazu na mechaniku newtonovskou. Ukazuje se však, že je principiálně nemožné formulovat základní koncepty kvantové mechaniky bez odkazu na mechaniku klasickou. Skutečnost, že elektron\footnote{V této a následujících sekcích budeme pro stručnost hovořit o ``elektronu'' ve smyslu jakéhokoliv objektu kvantové povahy, tj. kupř. částice nebo soustavy částic následující zákony kvantové, nikoli klasické, mechaniky.} nemá žádnou trajektorii v klasickém smyslu, implikuje taktéž, že sám o sobě nemá žádnou dynamickou charakteristiku.\footnote{Zde hovoříme o veličinách, které popisují pohyb elektronu, nikoli o veličinách vztažených k elektronu jako k částici, např. náboji nebo hmotnosti; tyto budeme označovat jako parametry.} Je tedy patrné, že pro systém složený pouze z kvantových objektů by bylo nemožné zkonstruovat jakoukoliv logicky nezávislou mechaniku. Konstrukce kvantitativního popisu pohybu elektronu tedy vyžaduje přítomnost fyzických objektů, jež dostatečně přesně následují zákony klasické mechaniky. Interaguje-li elektron s takovýmto ``klasickým objektem'', dochází obecně ke změně stavu tohoto klasického objektu. Charakter a velikost této změny potom závisí na stavu elektronu, lze ji tedy užít pro jeho kvantitativní popis.\par
V rámci tohoto vztahu obvykle ``klasický objekt'' nazýváme \emph{měřicím přístrojem} a jeho interakci s elektronem \emph{měřením}. Zde je nutné vyzdvihnout, že nemyslíme proces měření, v němž fyzik-pozorovatel hraje aktivní roli. \emph{Měřením} v kvantové mechanice rozumíme jednoduše jakoukoliv interakci mezi klasickými a kvantovými objekty, k níž dochází odděleně a nezávisle na jakémkoliv pozorovateli. Důležitost konceptu měření v kvantové mechanice první osvětlil N. Bohr.\par
Jakožto ``měřicí přístroj'' jsme definovali fyzický objekt, který s dostatečnou přesností následuje zákony klasické mechaniky. Tímto může ku příkladu být těleso s dostatečně velkou hmotností. Měřicí přístroj však nemusí nutně být makroskopickým. Za určitých okolností může roli měřicího přístroje hrát také objekt mikroskopický, neboť podmínka ``dostatečné přesnosti'' je závislá na konkrétním problému. Například pohyb elektronu ve Wilsonově komoře je pozorován prostřednictvím mlžné stopy, jejíž rozměry jsou velké ve srovnání s rozměry atomu; je-li dráha určena s takto malou přesností, lze elektron popsat plně klasicky.\par
Kvantová mechanika mechanika tedy mezi fyzikálními teoriemi zaujímá zvláštní postavení: klasická mechanika je jejím limitním případem, zároveň však tento limitní případ vyžaduje pro svou formulaci.\par
Nyní již můžeme formulovat úlohu kvantové mechaniky. Obvyklá úloha spočívá v předpovědi výsledků následujících měření prostřednictvím výsledků měření předchozích. Později navíc uvidíme, že ve srovnání s mechanikou klasickou, kvantová mechanika navíc omezuje rozmezí hodnot, kterých mohou různé fyzikální veličiny nabývat (např. energie): tj. hodnot, které lze získat prostřednictvím měření dané veličiny. Metody kvantové mechaniky nám tedy musí umožnit najít tyto povolené hodnoty.\par
Proces měření v kvantové mechanice má velmi významnou vlastnost: vždy ovlivní měřený elektron a je principiálně nemožné pro danou přesnost měření tento efekt učinit libovolně malým. Čím přesnější je měření, tím silnější je vliv na měřený objekt, pouze pro velmi nepřesná měření může být ovlivnění měřeného objektu malé. Tato vlastnost měření logicky souvisí se skutečností, že dynamické vlastnosti elektronu vyvstávají pouze jakožto výsledek měření sám. Je patrné, že pokud by bylo možné vliv měření na objekt učinit libovolně malým, byla by měřená veličina sama určitou hodnotu nezávislou na měření.\par
Mezi různými druhy měření můžeme vyzdvihnout roli měření souřadnic elektronu. V rámci kvantové mechaniky lze měření souřadnic elektronu vždy provést\footnote{Zde ještě jednou upozorňujeme, že ``provedením měření'' myslíme interakci elektronu s klasickým ``měřicím zařízením'', a nečiníme tedy předpoklad přítomnosti vnějšího pozorovatele.} s libovolnou přesností.\par
Předpokládejme, že po určitých časových intervalech \(\Delta t\) provádíme postupná měření souřadnic elektronu. Výsledky nebudou obecně ležet na hladké křivce. Naopak, čím přesněji budeme souřadnice měřit, tím neuspořádanější a nespojité budou naměřené výsledky, a to v souladu s neexistencí trajektorie elektronu. Hladkou dráhu získáváme pouze v případě, kdy jsou souřadnice elektronu měřeny nepřesně, jako tomu bylo v případě kondenzované páry ve Wilsonově komoře.\par
Pokud naopak zanecháme přesnost měření neměnnou a budeme zmenšovat intervaly \(\Delta t\) mezi měřeními, budou sousední měření samozřejmě odpovídat sousedním souřadnicím. Přestože však budou výsledky řady měření ležet v některé malé oblasti prostoru, jejich rozložení v této oblasti bude plně neuspořádané, nebude odpovídat žádné hladké křivce. Specificky, uvažujeme-li \(\Delta t\) blížící se nule, výsledky sousedních měření se neblíží žádné přímce.\par
Tato skutečnost ukazuje, že v kvantové mechanice neexistuje koncept rychlosti částice v klasickém slova smyslu, tj. limity, k níž směřuje rozdíl souřadnic naměřených ve dvou okamžicích, pokud časový interval \(\Delta t\) mezi těmito okamžiky směřuje k nule. Později uvidíme, že i přes tento problém lze v kvantové mechanice zkonstruovat rozumnou definici okamžité rychlosti částice, která s přechodem v klasickou mechaniku přechází v klasickou okamžitou rychlost. V klasické mechanice můžeme částici však v jednom okamžiku přiřadit určitou polohu a určitou okamžitou rychlost zároveň, kdežto v kvantové mechanice je situace zcela jiná. Pokud bychom měřením nalezli přesné souřadnice elektronu, nebylo by možné určit jeho rychlost. Pokud bychom naopak určili jeho rychlost, nebylo by možné určit jeho polohu v prostoru. Současná existence polohy a rychlosti by vedla k existenci trajektorie, kterou elektron nemá. V kvantové mechanice jsou tudíž souřadnice a rychlost elektronu hodnoty, které nelze přesně změřit současně, tj. nemohou mít současně přesně dané hodnoty. Můžeme říci, že souřadnice a rychlost elektronu jsou veličiny, jež současně neexistují. V následujícím potom odvodíme kvantitativní vztah, který popisuje možnost nepřesného současného měření polohy a rychlosti.\par
Úplný popis stavu fyzikálního systému v klasické mechanice je proveden určením všech souřadnic a rychlostí v daném okamžiku; prostřednictvím těchto počátečních podmínek můžeme pohybovými rovnicemi popsat chování systému ve všech následujících okamžicích. V kvantové mechanice je takovýto popis v principu nemožný, neboť souřadnice a odpovídající okamžité rychlosti nemohou existovat současně. Popis stavu kvantového systému je tedy nutně proveden menším počtem veličin než v mechanice klasické, tj. je méně detailní než klasický popis.\par
Z tohoto vyplývá velmi důležitý důsledek o povaze předpovědí v kvantové mechanice. Zatímco klasický popis dostačuje pro předpověď budoucího pohybu mechanického systému s dokonalou přesností, méně detailní kvantový popis pro tento účel evidentně nemůže postačovat. V důsledku, třebaže je učiněn nejkompletnější možný popis stavu elektronu v kvantové mechanice, nelze s jistotou předpovídat jeho chování v následujících okamžicích. Kvantová mechanika tudíž nedokáže činit přesné předpovědi budoucího chování elektronu. Pro daný počáteční jeho stav mohou následná měření dávat různé výsledky. Úloha kvantové mechaniky potom spočívá v určení pravděpodobností různých výsledků těchto měření. Poznamenejme, že v některých případech může samozřejmě pravděpodobnost určitého výsledku být rovna jedné, tj. s jistotou obdržíme určitý výsledek.\par
Veškerá měření v kvantové mechanice můžeme rozdělit do dvou tříd. V první třídě, která obsahuje většinu měření, nacházíme taková, v nichž žádný stav systému nevede s jistotou k určitému výsledku. Druhá třída potom obsahuje měření, kde pro každý výsledek měření existuje stav systému, který k tomuto výsledku vede s jistotou. Měření z druhé třídy označujeme jako \emph{predikovatelná}, hrají v kvantové mechanice důležitou roli. Kvantitativní charakteristiky stavu určené takovýmto měřením pak v kvantové mechanice nazýváme \emph{fyzikálními veličinami}. Pokud v některém stavu dává měření s jistotou některý výsledek, tvrdíme, že má odpovídající veličina v tomto stavu danou hodnotu. V dalším budeme vždy termín ``fyzikální veličina'' užívat ve smyslu uvedeném zde.\par
Budeme se často setkávat se skutečností, že ne každá sada fyzikálních veličin v kvantové mechanice lze měřit současně, tj. všechny veličiny nemohou nabývat jednoznačných hodnot zároveň. Již jsme zmínili jeden příklad, a to rychlost a polohu elektronu. Důležitou roli v kvantové mechanice hrají sady veličin s následující vlastností: tyto veličiny lze měřit současně, avšak určíme-li současně jejich hodnoty, žádná jiná veličina (kromě funkcí veličin z této sady) nemůže mít v dané sadě jednoznačnou hodnotu. Tyto sady veličin nazýváme \emph{úplnými}.\par
Jakýkoliv popis stavu elektronu získáváme jakožto výsledek nějakého měření. Nyní zformulujeme definici \emph{úplného popisu} stavu v kvantové mechanice. Úplně popsané stavy jsou výsledky současného měření úplné sady fyzikálních veličin. Z výsledků takového měření lze zejména určit pravděpodobnosti výsledků jakýchkoliv následných měření, a to nezávisle na stavu elektronu před prvním měřením. Dále (s výjimkou §14) budeme stavy kvantového systému rozumět pouze tyto úplně popsané stavy.
\section{Princip superpozice}
